% ---
% title: A shooting contest
% tags: independence
% ---

\begin{problem}
	Alice and Bob shoot at a target in turns, and the first one to hit the
	target wins. Alice hits with probability $1/2$, Bob hits with probability
	$1/3$, and every shot is independent of all the other shots. Since Bob is
	the weaker shooter, Alice lets him shoot first. So the order of the shots
	is Bob, Alice, Bob, Alice, and so on.
	(a) What is the probability that Bob wins?
	(b) Is the contest fair, that is, do both players win with the same
	probability?
\end{problem}

\begin{solution}
	\idea Split the event ``Bob wins'' by the shot on which he wins. These
	events are mutually exclusive, so their probabilities add up. Each of them
	is an intersection of independent shots, so its probability is a product.
	The sum is a geometric series.
	\step{(a)} Call one \emph{round} a shot of Bob followed by a shot of
	Alice. Both players miss in a round with probability
	\[ P(\text{Bob misses})\cdot P(\text{Alice misses})
		=\frac23\cdot\frac12=\frac13, \]
	because the two shots are independent.

	Bob wins with his first shot if he hits it: probability $1/3$.

	Bob wins with his second shot if he misses, Alice misses, and then Bob
	hits. The three shots are independent, so we multiply:
	\[ \frac23\cdot\frac12\cdot\frac13=\frac13\cdot\frac13. \]

	In general, Bob wins with his $(k+1)$-th shot ($k=0,1,2,\dots$) if the
	first $k$ rounds are complete misses and then Bob hits. By independence
	this has probability
	\[ \Bigl(\frac23\cdot\frac12\Bigr)^{k}\cdot\frac13
		=\Bigl(\frac13\Bigr)^{k}\cdot\frac13. \]

	``Bob wins with his first shot'', ``Bob wins with his second shot'', and so
	on cannot happen together (the contest stops at the first hit), so these
	events are mutually exclusive. The probability of their union is the sum
	of their probabilities. This is a geometric series with first term $1/3$
	and ratio $q=1/3$. Since $|q|<1$, its sum is (first term)$/(1-q)$:
	\[ P(\text{Bob wins})
		=\sum_{k=0}^{\infty}\Bigl(\frac13\Bigr)^{k}\cdot\frac13
		=\frac{1/3}{1-1/3}=\frac{1/3}{2/3}=\frac12. \]
	\step{(b)} Someone wins with probability $1$: the probability that
	nobody hits in the first $n$ rounds is $(1/3)^n$, which goes to $0$.
	So Alice wins exactly when Bob does not win:
	\[ P(\text{Alice wins})=1-P(\text{Bob wins})=1-\frac12=\frac12. \]
	Both win with probability $1/2$: the contest is fair.

	\emph{A shortcut.} Look at one round only. Inside a round,
	\[ P(\text{Bob hits first})=\frac13,\qquad
		P(\text{Bob misses, then Alice hits})=\frac23\cdot\frac12=\frac13. \]
	These two probabilities are equal. If nobody hits, the next round starts
	in exactly the same situation, because the shots are independent. So in
	whichever round the contest is decided, Bob and Alice are equally likely
	to be the winner. Hence each wins with probability $1/2$.
	\begin{answer}
		(a) $P(\text{Bob wins})=1/2$ \quad (b) yes, it is fair:
		$P(\text{Alice wins})=1/2$ as well
	\end{answer}
	\begin{remark}
		\textbf{Common mistake:} comparing only the hit probabilities
		($1/3<1/2$) and concluding that Alice is the favorite. Who
		shoots first matters too. If Alice shot first, the same method would
		give
		\[ P(\text{Alice wins})=\frac{1/2}{1-1/3}=\frac34, \]
		so she would be a clear favorite. Letting Bob start balances the game
		exactly.
	\end{remark}
\end{solution}
